% ============================================================
%  K. Kroman, "A Few Further Remarks on 'The Principle of Relativity'"
%  Fysisk Tidsskrift, 15. Aargang (1916-17), pp. 192-205.
%
%  ENGLISH TRANSLATION
%  Status: COMPLETE. Translated from transcription.tex (the Danish diplomatic
%  transcription, COMPLETE), never from the scan. Structure mirrors
%  transcription.tex 1:1: same paragraphing, same page markers at the same
%  joints, the five footnotes at the same anchors with the same marks and the
%  same inconsistent placement about the stop.
%
%  Kroman's reply to the three articles that answered his 1915 attack --
%  Holst's »Tidsproblemet«, Walsøe's »Omkring Videnskabens Grænse« and
%  H. M. Hansen's »Relativitetsprincipet« -- and the close of the debate.
%  The first two are translated in ../holst-tidsproblemet and
%  ../walsoe-videnskabens-graense; Hansen's remains under the repo's embargo
%  until 1 January 2027, so this reply can be read against two of its three
%  targets only. See ../relativitetsprincipet/note.md.
%
%  RIGHTS: clear. Kroman died 26 July 1925; Danish copyright expired at the
%  end of 1995.
%
% ------------------------------------------------------------
%  TRANSLATION CONVENTIONS -- this article
% ------------------------------------------------------------
%  * Standing method: ../../../TRANSLATION-PLAYBOOK.md; the cluster's shared
%    decisions are set out at length in
%    ../relativitetsprincipet/translation.tex.
%  * EMPHASIS. All 42 \emph{} runs carried over, including the one that sits
%    INSIDE a compound on p. 195 -- Bevægelses\emph{ændrings}aarsag, rendered
%    "a cause of \emph{change} of motion" so that the emphasis still falls on
%    the same element -- and the two on p. 205 that enclose the italic
%    variables A and B.
%    NOTE: unlike the 1915 article, this one letterspaces NO proper names at
%    all on pp. 192-193 or 200, and the compositor sets »Einstein'ske« there
%    in ordinary roman. The English carries no \emph{} on those names either;
%    that is not an oversight.
%  * "Dr. H." is Kroman's abbreviation for Dr. H. M. Hansen throughout, "L."
%    for Lorentz and "M." for Michelson. All kept as he sets them.
%  * "Fænomenist" -> "phenomenist", the term Kroman is attacking, kept rather
%    than modernised to "phenomenalist"; his own gloss follows in the text.
%  * The GERMAN quotations (pp. 198, 202) stay in German, inside quotation
%    marks and upright, as printed.
%  * The old Danish id-est sign ɔ: (pp. 195, 202) -> "i.e." per the playbook.
%  * PRINTER'S DEFECTS. p. 203's »sikkert!,« (exclamation mark and comma
%    together) and p. 204's misplaced comma before "that" are reproducible and
%    are reproduced. p. 200's spaced product »2 l a« against p. 201's »2la« is
%    kept as printed in both places.
%  * FOOTNOTE MARKS sit inconsistently about the full stop -- BEFORE it on
%    pp. 193 and for the **) on p. 198, AFTER it for the *) on p. 198 and on
%    p. 205. Reproduced in each place.
%
%  BUILD: pdflatex (libertinus + libertinust1math), like the Danish.
% ============================================================

\documentclass[12pt,a4paper]{article}
\usepackage[utf8]{inputenc}
\usepackage[T1]{fontenc}
\usepackage{libertinus}
\usepackage{libertinust1math}
\usepackage{textalpha}
\usepackage{amsmath}

\usepackage{graphicx}    % for \reflectbox, below
% The old Danish id-est sign »ɔ:« is rendered "i.e." in the English per the
% playbook, so the declaration is not needed for the body; kept so the file
% compiles unchanged if a later editor restores the sign in a quotation.
\DeclareUnicodeCharacter{0254}{\reflectbox{c}}

% FOOTNOTES. Marked *) on pp. 193, 198, 200, 205. p. 198 carries a second,
% marked **); that one is given its mark inside a group at the site rather
% than by resetting the counter, which keeps hyperref's anchors unique.
\renewcommand{\thefootnote}{*)}
\usepackage[english]{babel}
\usepackage{enumitem}
\usepackage[protrusion=true,expansion=true]{microtype}
\usepackage{setspace}
\usepackage[margin=1.2in]{geometry}
\usepackage{fancyhdr}
\setlength{\headheight}{15pt}
\usepackage{hyperref}

\hypersetup{
  pdftitle={A Few Further Remarks on The Principle of Relativity},
  pdfauthor={K. Kroman},
  hidelinks
}

\pagestyle{fancy}
\fancyhf{}
\fancyhead[LE,RO]{\thepage}
\fancyhead[RE]{\textit{A Few Further Remarks}}
\fancyhead[LO]{\textit{K. Kroman}}

\setstretch{1.3}

\title{\textbf{\textbf{A Few Further Remarks on}\\[2pt]``The Principle of Relativity''}}
\author{K.\ Kroman}
\date{\textit{Fysisk Tidsskrift} 15 (1916--17), pp.~192--205\\[2pt]
\small Translated from the Danish}

\begin{document}
\maketitle

% --- p. 192 ---
% SHARED PAGE. The upper two thirds of printed p. 192 belong to the previous
% article, P. O. Pedersen, »Den elektriske Bues Elektronteori«. Kroman's
% article begins about three fifths of the way down, with its own centred
% title block, which \maketitle supplies. Nothing of Pedersen is translated.
Fysisk Tidsskrift opens its 15th volume with no fewer than three treatises on
the \emph{Einstein}ian principle of relativity. Since all three of them take
up a position towards my article on the same subject in the 14th volume,
pp.~1---30, I should like to be allowed to bring forward a few further remarks
on the question.
% *** DISCREPANCY IN transcription.tex, FOR THE EDITOR TO SETTLE ON THE IMAGE.
% The Danish BODY sets »det \emph{Einstein}'ske Relativitetsprincip« -- one
% letterspaced run -- while the comment immediately below it in the same file
% says the opposite: "no letterspacing anywhere on p. 192 or p. 193 ...
% It is transcribed WITHOUT \emph." The header's emphasis note ("42 runs, and
% none at all on pp. 192-193 or 200") agrees with the COMMENT, not with the
% body; the file's own run count of 42 agrees with the BODY.
% This translation mirrors the BODY, transcription.tex being the source of
% truth, so the run is carried here. If the image shows ordinary roman, drop
% the \emph{} in BOTH files and the total falls to 41.

With the author of the first treatise, Mr.\ Editor Helge Holst, I am, I
believe, in agreement on each and every point. His introductory remark about
the practical applicability of the Einsteinian formulae I can also readily
subscribe to; indeed I have said the same thing myself at the top of p.\ 24 of
my article.

Nor am I, with the author of the second treatise, Mr.\ Engineer
C.\ E.\ Walsøe, at bottom in disagreement so far as the Einstein problem is
concerned. When he (at the foot of p.\ 15) finds that Einstein's doctrine has
been condemned as a scientific hypothesis, I am completely satisfied. To
calling it an unscientific hypothesis I have no objection.

% --- p. 193 ---
With the author of the third treatise, Dr.\ H.\ M.\ Hansen, I am on the other
hand, epistemologically regarded, in great disagreement. Dr.\ H. appears as a
warm defender of Einstein's doctrine, which he does not find very different
from Lorentz's, and he holds that I have misunderstood the latter. I assume
that this is chiefly to mean that I have not sufficiently distinguished
between Lorentz's original and his later, more fully worked out hypothesis.
But to this my article gave me no occasion whatever. Dr.\ H. overlooks that I
expressly emphasised that I did not wish to write a theory of electricity, but
merely to examine the Einsteinian hypothesis, particularly on its
epistemological side. Neither did I wish to give an exposition of Lorentz's
works, the less so as I am not decidedly certain that they give us precisely
the right solution of the Michelson riddle\footnote{It really continues to be
my hope that M.'s result will one day be explained by an alteration of the
theory of light. Dr.\ H. finds (p.\ 18) that this prospect is barred by an
utterance of Poincaré's. To me, however, it seems that this utterance --- at
least in Dr.\ H.'s formulation --- is quite unwarranted. By what right could
one demand that L.'s hypothesis should take account of other marvels than
M.'s, before any others had yet shown themselves? If others of the same nature
turn up, the same hypothesis will of course be able to explain them too, and
if they are of a new nature, new hypotheses will presumably be required as
well. But one should not, as is well known, sell the skin before one has shot
the bear.}.
% FOOTNOTE MARK POSITION: the printing sets »Gaade*).« -- the mark BEFORE the
% full stop. Compare pp. 198 and 205, where it stands after it. Followed.

Lorentz has presented both his hypotheses --- the latter no fewer than three
times --- in the writings I have cited. He has there discussed the various
experiments akin to Michelson's which Dr.\ H. also mentions. I can therefore
hardly have overlooked that he had both an original and a more fully worked
out hypothesis. But common to them both is the contraction, and I needed
nothing else and nothing more for my examination of Einstein's train of
thought. I have moreover also emphasised the community of formulae between
Einstein and Lorentz, but at the same time stressed that in spite of this
there is nevertheless a most essential difference between the two. In what
follows I shall once more try to give grounds for holding that Lorentz's
hypotheses are at all events possible assumptions, while Einstein's must in my
view be stamped as impossible.

Dr.\ H. calls Einstein's doctrine simpler and (in an earlier treatise) bolder
than L.'s. That the simplicity must be of a rather
% The foot of p. 193 carries the signature line »Fysisk Tidsskrift. XV« with
% the sheet number »14«; p. 195 carries »14*«. Not transcribed.

% --- p. 194 ---
superficial kind seems, however, to appear sufficiently from the uncommonly
artificial results concerning space and time in which it issues, and boldness
is after all no unmixed scientific virtue. In Einstein it appears, as we shall
see, chiefly as courage for exaggerations and vague determinations of
concepts, courage to defy all reasonable probability, and indifference towards
fundamental epistemological demands. I have, in my judgement, brought this
last point to light in a number of places in my treatise. Dr.\ H. has sought
to remove some of the difficulties pointed out; but it does not seem to me
that he has really succeeded in doing so. I shall however not dwell further on
these particulars, since I believe I can illuminate and justify my view still
more clearly by following Einstein's train of thought critically through its
development, and pointing out where its incorrectnesses come forward one by
one and give occasion to the various contradictions and great improbabilities
in which his doctrine issues.

First, however, I must permit myself a couple of general remarks.

Dr.\ H. defends various peculiarities in Einstein's exposition with the
statement that Einstein is a phenomenist. That is however a little reassuring
defence; for it is anything but fortunate for an investigator to be a
phenomenist. By a phenomenist one understands an investigator who has formed
the view that it is most scientific to keep to the phenomena, to that ``which
shows itself'', and to let causes, forces and the like look after themselves.
It is thus sensation that is hereby set in the high seat, while thought, as an
unreliable counsellor, is put in the corner. But of course such an
epistemology is just as impossible to carry through as it is false. Without
the help of thought there could of course be neither calculation nor
inference. A real denial of thought therefore becomes an impossibility for
man, and every phenomenist therefore becomes of necessity an
\emph{inconsistent} investigator, who lets thought play its part as usual the
99 times, while suddenly the 100th time it is quite arbitrarily denied, so
that one permits oneself, for example, to set up effects without stating or
being able to state a trace of causes, and the like. Of this kind of step we
shall meet a number in Einstein, and that such liberties are not fortunate for
scientific character surely needs no demonstration.

% --- p. 195 ---
Let us then begin by looking a little more closely at the not very clear
concept and the not very fortunate expression ``the principle of relativity'',
and let us first ask: did Newton (and the older physics generally) have
anything that could rightly and reasonably be called a principle of
relativity? What, in general, is a principle?

In most physical textbooks one finds various quite respectable propositions
obstinately entitled principles: ``Archimedes' principle'', ``Doppler's
principle'', and so on. This evidently arises solely from linguistic
indifference and unfortunate inertia. By a principle is understood, more
exactly speaking, a starting point, a fundamental proposition, i.e. a
proposition which is laid at the foundation of the whole following structure
of doctrine, a first proposition which is not derived from earlier ones, which
has no proof before it, since one has found it self-evidently correct, and of
whose validity no one has dared to doubt. Of such principles or fundamental
propositions Newton, as is well known, set up three: the law of inertia, the
law of force and the law of interaction. A ``fundamental proposition of
relativity'', on the other hand, he did not set up. In the law of force, and
already earlier, he defined a force not as a cause of motion, but as a cause
of \emph{change} of motion, thus not by the aid of motion but by the aid of
% EMPHASIS: the letterspaced run is »æ n d r i n g s« ONLY, standing inside
% the compound »Bevægelses æ n d r i n g s aarsag«. The English keeps the
% emphasis on the same element: "a cause of \emph{change} of motion".
change of motion, that is, change of velocity, of direction, or of both. The
logical consequence of this is, among other things, that if a body has a
certain motion, and another like it has the same motion plus some uniform
rectilinear velocity, the equations of force for both must of course become
the same; for unaltered motion requires no force.

Further, he maintains that a body can have motions of all kinds, relative as
well as absolute, and that just as there is a difference between uniform
rectilinear motion with the velocity 5 and uniform rectilinear motion with the
velocity 10, just so must ``motion with the velocity 0'' be something by
itself. This too every sober person is compelled to assume. No one can think
anything by the expression motion without at the same time having to think,
more or less distinctly, of a resting; and no one can think a motion as
relative without at the same time having to think an absolute rest or motion.
These ideas stand and fall with one another like the ideas up and down, and in
every rational mechanics one handles these determinations quite securely and
naturally.

% --- p. 196 ---
\emph{Widely different from this} is however the question: can we human beings
always determine what motion a body has in a given real case? With absolute
exactness we can of course never do it. With approximation, on the other hand,
we can often determine its relative motion. In certain cases we shall also be
able to determine that it has such and such an absolute motion. But if the
motion is uniform and rectilinear, the case will --- so far as we yet know ---
not be distinguished by any sensible mark from an unaltered motion in the same
direction with another velocity or with none at all, and therefore it will, as
Newton says, often be \emph{hard} to determine the true velocity.

Now why may Newton have used so vague an expression? Why has he not said
outright that in the face of unaltered motion a determination of velocity is
simply and plainly impossible? Of course he was quite clear that here no help
was to be expected from manifestations of inertia, as there is with changes of
velocity or of direction. He had, it is true, the \emph{conjecture} that the
centre of gravity of the solar system was absolutely at rest in space. But it
was only a conjecture. He has also said that perhaps there was no such fixed
point at all that we could lean upon. He has therefore, rather, presumably
thought that matter, which had so lately revealed itself from a new side in
the enigmatic gravitation, might perhaps also one fine day disclose
\emph{certain} differences according as it was at rest or in uniform
rectilinear motion, and that the physicists might learn to control these
differences, so that a computation of velocity would thereby become possible.
If Lorentz's hypothesis is correct, we have of course in \emph{the contraction}
precisely something in the direction indicated.

When later the wave theory of light had come to power, and the phenomenon of
interference had been more closely understood, a number of physicists have
surely thought that as soon as one merely obtained sufficient technical
mastery of the interference experiments, one would thereby have a means of
distinguishing between absolute rest and uniform rectilinear motion. Precisely
this was of course Michelson's assumption.

If one wishes to express oneself as exactly as possible, I therefore think
that one must confine oneself to saying of Newton and the older physicists:

1) They conceded that they were for the present unable \emph{experimentally}
to distinguish between absolute rest and uniform rectilinear motion.
% The spaced run »e k s p e r i-/m e n t e l t« is broken over a line turn in
% the printing; it is one run, and is one run here.

% --- p. 197 ---
2) But they were by no means certain that this would remain impossible.

3) Still less did they set up the impossibility of it as a principle, that is,
as a proposition one \emph{was to} start from.

4) And they have surely all been agreed with Newton that absolute rest and
uniform rectilinear motion were something different in themselves, and not
simply dependent upon which coordinate systems one used.

But if this is how the matter stands, it becomes not merely misleading but
quite incorrect to use expressions like ``a Newtonian principle of
relativity''. Everything is then said far more clearly and sharply in the
simple words that Newton built upon the law of inertia. ---

More characteristic than to call Newton a relativist would surely be to call
him an absolutist. In Einstein's doctrine, on the other hand, we have nothing
less than an attempt to relativise not merely rest and motion, but even space
and time as well.

It is in particular Michelson's result that leads him into these endeavours.
Michelson started from the assumption that light, independently of the light
source's rest or motion, constantly had in the world-space and in air the
constant velocity $c$, as he relied on the current theory of light generally.
So did Lorentz and Einstein, and I can therefore frame the following remarks
on the basis of the inviolability of this presupposition.

Michelson now found by his experiment that light traversed the two paths $s$
and $s_{1}$ in one and the same time.

For human reason there is nothing else to remark upon this than that the two
paths consequently \emph{must} have been equally long. This was also
Michelson's first thought, in that he assumed that the earth's translation had
been counteracted by the motion of the whole solar system. When this way out
was found barred, Lorentz set up his contraction hypothesis. If we leave aside
that he might of course also have assumed the shorter path to be lengthened,
and if we hold fast to the complete correctness of the theory of light, there
was, as has been said, no other way out, unless the phenomenon was to come
under the purely irrational.

Einstein has surely known this first contraction hypothesis of Lorentz's when
in 1905 he came forward with his doctrine. This rests, he says at the
beginning of his treatise, upon two presuppositions: the principle of
relativity and the ``only apparently irreconcilable''

% --- p. 198 ---
presupposition of the aforementioned constant velocity of light. These two
presuppositions are \emph{sufficient}, he adds. Of contraction he says not a
word. On the other hand he declares it superfluous to assume an
ether.\footnote{Lorentz, Einstein, Minkowski: Das Relativitätsprincip,
pp.\ 27---28.}
% FOOTNOTE MARK POSITION: printed »Æter.*)« -- mark AFTER the full stop, where
% p. 193 sets it before. Both followed as printed.

Concerning ``the principle of relativity'' he remarks that he has arrived at
the conjecture that to the concept of absolute rest there corresponds, neither
in the mechanical nor in the electrodynamic domain, any peculiarity of the
phenomena. This conjecture he will now ``raise to a principle'' and call the
principle of relativity.

But it really will not do to raise conjectures to ``principles'' in this way
without further ado. For the present we know only that in a few cases results
corresponding to Michelson's have shown themselves. Einstein can therefore
only be permitted to set up as a \emph{hypothesis} that his conjecture holds
generally, and he can thereupon set himself the task of testing this
hypothesis's possibility and its eventual scientific worth. But he thereby
incurs quite other obligations than if he could at once raise his conjecture
to a principle.

His first scientific obligation would thus become that of giving the
hypothesis up again on account of its impossibility. For it is impossible,
since it stands fast that, when the correctness of the theory of light is
presupposed, the contraction is the only conceivable way out. But if bodies
are contracted by the motion, then there corresponds after all a definite
peculiarity, namely uncontractedness, to real rest. That we cannot for the
present measure the contraction \emph{directly} is, practically regarded,
tiresome enough. Scientifically regarded, however, it is already sufficient
that we can establish that it must take place, indeed even compute its
magnitude in relation to any presupposed velocity.

But with this, as we shall see, all real relativising of rest and motion, of
space and time, is annihilated, and the many corresponding exaggerated
expressions are to be rejected{\renewcommand{\thefootnote}{**)}\footnote{Dr.\ H.
too uses these misleading expressions. Repeatedly it is said: rest and uniform
rectilinear motion are one and the same, or physically one and the same; there
is no sense in distinguishing between them. And yet he then distinguishes so
vigorously between them that he maintains that the motion makes bodies
contract, clocks go more slowly, people age more slowly; indeed he even adds
(p.\ 22) that every observer can establish this.}}. Already here, then, it is
seen that Einstein's hypothesis is an impossibility. What he could with
scientific
% THE ONLY PAGE IN THE ARTICLE WITH TWO FOOTNOTES, marked *) and **).
% \thefootnote is fixed at *) in the preamble, so the second is given its mark
% inside a group here rather than by resetting the counter, which keeps
% hyperref's anchors unique. Printed »forkastelige**).« -- mark BEFORE the
% full stop, the opposite of »Æter.*)« above.

% --- p. 199 ---
right have set up and sought to justify is something widely different from
what he has stated. He could rightly have set up the hypothesis that there is
a certain \emph{technical} equivalence between absolute rest and uniform
rectilinear motion.

Without noticing anything of the contradiction, Einstein however goes further.
He thinks of two coordinate systems with mutually parallel axes and a common
$X$-axis, the one $(O\,x\,y\,z)$ at rest, the other $(O'\,x'\,y'\,z')$ with
the constant $X$-velocity $v$. In each there is an observer ($A$ and $B$)
furnished with perfectly alike clocks and measuring rods and a consummate
ability to use them. They both rely on the theory of light and on the constant
velocity of light $c$. Dr.\ H. calls them physicists and assumes that they
know Michelson's experiment. All this we can presuppose without misgiving.
% MATHEMATICS: the printing sets the coordinate letters in italic with a thin
% space between them and the parentheses in roman -- (O x y z), (O' x' y' z').
% Reproduced with \, here as in the Danish.

But Einstein then continues: they both assume that they are at rest.

This seems a remarkably frivolous and dogmatic assumption to assign to them.
Physicists so excellently equipped ought surely rather to have repeated
Michelson's experiment critically, and in any case to have built upon the more
reasonable presupposition that each of them had a uniform rectilinear motion
with a certain unknown velocity (perhaps $= 0$ for one of them), which they
would have either to keep as an unknown in their calculation or to try to
determine as well as could be done, for example by the aid of a sky of fixed
stars, or whatever else they had that was most serviceable.

We may, however, suppose that Einstein has assigned them these arbitrary
assumptions with a definite end in view, perhaps in order to show quite
clearly that rest and uniform rectilinear motion are entirely alike, or the
like; and of course he is at liberty in his thought-experiment to introduce
whatever determinations he pleases, provided only he informs the reader what
he is doing.

But at the same time the reader is of course also at liberty to note what
happens: at liberty to note that here Einstein introduced a contradiction into
the situation, in that he gave $A$ a correct presupposition and $B$ a false
one, in direct contradiction with what had been said about the two coordinate
systems.

$A$ thereupon regulates the clocks he has set up round about in the system $O$
by the aid of light signals, so that they all get the same setting and a
correct rate, 1 second while the light runs the path $c$. This now goes easily
enough. But then $B$ is to do the same from his contradiction-laden
presupposition.

% --- p. 200 ---
And here one must remember that Einstein has hitherto not said a word about
contraction, as indeed he could hardly do, since he has on the contrary said
that the phenomena were to behave alike under rest and under uniform
rectilinear motion. We must therefore for the present reckon without any
contraction. Einstein has of course possibly thought that it would present
itself of its own accord as a consequence of his presuppositions.

But suppose now, first, that the ``perfect'' clocks have also had a correct
rate (1 second while the light really runs the path $c$) already from the
manufacturer's hand, and that $B$ has relied on this as well as on the
constant velocity of light. What will then happen?

$B$ would then simply discover that if he sent a light signal the path $l$
there and back perpendicular to the $X'$-axis, the light would not be
$\dfrac{2l}{c}$ seconds on the way, but $\dfrac{2l}{c}\,a$
seconds\footnote{Compare, concerning this and the following expressions, my
earlier article.}; and if he sent the signal the path $l$ there and back along
the $X'$-axis, it would take the time $\dfrac{2l}{c}\,a^{2}$. But from this
$B$ would surely conclude: my system therefore cannot be at rest; it must have
the uniform rectilinear velocity $v$.

Assuming next that $B$ does not rely on his clocks' going correctly, but will
set them himself, we get a somewhat altered result. Let him again begin
perpendicular to the $X'$-axis and again use the measured path $l$. If he now
arranges his clock to advance $\dfrac{2l}{c}$ second-marks while the light is
on its way, his seconds will be $a$ times too large; for the light has in
reality then run the path $2\,l\,a$. His clock thus comes to go $a$ times too
% The printing sets this product with visible spaces, »2 l a«, where p. 201
% sets the same quantity closed up as »2la«. Both reproduced as printed.
slowly, of which however he for the present has no inkling. But if he then
sends a signal along the $X'$-axis, he will to his surprise find that the
regulated clock has now advanced $\dfrac{2l}{c}\,a$ ``seconds'' while the
light was on its way. And if he has thought the matter over a little, he will
surely also conclude from this that he not only has the aforementioned
velocity $v$, but that he has moreover arranged his clock to go $a$ times too
slowly.

Thus must the results shape themselves when we reckon without contraction.

% --- p. 201 ---
But these results Einstein cannot use; for they stand in conflict with the
Michelson result.

\emph{Could} he introduce a real contraction, $B$'s self-deception could on
the other hand succeed. If all $B$'s distances in the $X'$-direction became
$a$ times as small, he would be able to arrange all his clocks to advance
$\dfrac{2l}{c}$ second-marks while the light, on his presupposition, ran the
path $2l$, but in reality the path $2la$. His clocks would thereby come to go
$a$ times too slowly, and his incorrect presupposition would further entail
that they got local time, or, in relation to the clock at $O'$, deficits of
clock-reading proportional to their $X'$-distances from $O'$, without his
noticing anything. All this is easily obtained by quite elementary
calculation, and in the same way one easily obtains what $A$ and $B$ would on
the whole be led to think about each other's computations, holding fast as
they do to their mutually conflicting presuppositions.

In this way, then, we should be able to arrive at ``the Einsteinian
transformation formulae''.

But Einstein himself cannot (legitimately) arrive at them. For he can neither
\emph{tolerate} nor \emph{dispense with} nor \emph{introduce} the contraction.

He cannot \emph{tolerate} the contraction; for it would annihilate his first
presupposition, which stands in conflict with it, since that presupposition
does not permit a difference between the phenomena under absolute rest and
under uniform rectilinear motion.

But even if Einstein could also \emph{tolerate} the contraction, he still
cannot \emph{introduce} it. For where should he get it from? Since he has
rejected the ether, he would be quite unable to find any cause for this
alteration of the phenomena. The contraction would become a purely miraculous
phenomenon; and surely one is not permitted in science to invent miracles in
order to prop up an otherwise tottering or unworkable hypothesis?

On the other hand we have just seen that the doctrine cannot at all
\emph{dispense with} the contraction. It is quite necessary if $B$ is to be
able in the long run to remain in his false assumption. And if this does not
happen, Einstein will not merely come into conflict with Michelson; his train
of thought will positively come to establish the opposite of what he set out
to prove.

% --- p. 202 ---
There is therefore hardly anything else to be done than to give up the
Einsteinian doctrine as \emph{impossible}.

Perhaps one or another of his adherents will remark that if there is nothing
more amiss than this, then the matter does not stand so badly for Einstein.
For we need then merely restrict his relativity presupposition a little, let
him keep the ether, and let him add the contraction; then all is in order.

Quite so, certainly. But then it is simply no longer Einstein's doctrine; then
it is Lorentz's that we have arrived at. And Lorentz's view is, as has been
said, a possibility, but in return, as we shall see, a view very different
indeed from Einstein's.

The difficulty with the contraction Einstein has evidently not had clearly in
view at all. When he comes to compute $B$'s \emph{moved} system, it goes quite
easily for him, ``indem man das Prinzip der Konstanz der Lichtgeschwindigkeit
im ruhenden Systeme anwendet.'' But that he has of course no right whatever to
% The German quotation is set upright, inside guillemets like every other
% quotation in the article -- NOT in italic. Followed.
do. He thereby commits a circular inference and proves nothing at all, or at
most that rest and motion are one and the same, provided one boldly
presupposes that they are one and the same. The error hangs together with his
``raising'' to a principle, without further ado, what is only a hypothesis. If
light, according to his second presupposition, has a constant velocity in
space and thus in the system at rest, then it is \emph{immediately regarded
unthinkable} that it should also have a constant velocity, i.e. the same
velocity in all directions, in a moved system. Anything of the kind first
becomes thinkable if the moved system and all its contents have, on account of
the motion, contracted according to certain definite laws. Such an alteration,
however, Einstein cannot get introduced without having to give up his first
presupposition and to assume the ether again. His two presuppositions are thus
\emph{really} in contradiction with each other, and not merely apparently so,
as he supposed.

Einstein could with every scientific warrant have introduced his problem with
the following question:

Can one, without coming into contradiction with sound reason, assume 1) that
the phenomena are always alike under absolute rest and uniform rectilinear
motion, and 2) that light always has in space (and in air) the same constant
velocity $c$?

And to this the answer would then have to be: no! For if the phenomena are to
be \emph{even merely apparently} alike under absolute rest

% --- p. 203 ---
and uniform rectilinear motion, then light must also in a uniformly
rectilinearly \emph{moved} system apparently have the velocity $c$ in all
directions. But this is unthinkable unless the system and all its contents are
contracted during the motion according to certain laws (for example the
Lorentzian ones). And if this is the case, then the phenomena are of course
not alike under absolute rest and uniform rectilinear motion. If it should
still be objected that the contraction cannot be reckoned to ``the
phenomenal'', to that which ``shows itself'', then one forgets that it has of
course shown itself precisely through the position of Michelson's interference
fringes.

But Einstein sees nothing of all this. He therefore keeps his presuppositions
without suspecting anything of their contradiction. He raises his impossible
hypothesis to a principle and takes the contraction along on his own free
authority. If he can assume the theory of light without assuming an ether, it
would of course also be unreasonable to make a fuss about such a trifle of a
contraction; if one is a phenomenist, such things cause no stoppages. And thus
he arrives at the result that he has really established that rest and uniform
rectilinear motion are, both in the mechanical and in the electrodynamic
domain, one and the same. For he has of course ``proved'' that $A$ and $B$
both arrive at the same results concerning the phenomena, indeed concerning
each other, although $A$ starts from a presupposition correct by the old
belief and $B$ from one completely wrong by the old belief. Consequently it
must after all be quite indifferent whether one presupposes oneself to be at
rest or in uniform rectilinear motion, and one can thus simply say to
Michelson: you need only presuppose that the earth stands still, and your
result at once becomes explicable. If anyone says: yes, but the earth
certainly does move!, you answer merely: rest and uniform rectilinear motion
% PRINTER'S USAGE, transcribed as printed: »ganske sikkert!, svarer Du blot«
% -- an exclamation mark followed by a comma. Reproduced.
% Note also that this whole reported dialogue is set WITHOUT quotation marks,
% where the same author quotes Einstein in guillemets elsewhere. Followed.
are quite the same. There is no sense in distinguishing between them.

This is really Einstein's explanation of Michelson.

But if rest and uniform rectilinear motion are thus one and the same, then
there is of course no sense either in calling $A$'s presupposition correct and
$B$'s false. Then they are both correct, and the contradiction between them is
something that will surely disappear with a more correct conception of the
nature of space and time. And it is thus not correct either that $B$'s seconds
are false. They are just as correct as $A$'s, although they are $a$ times as
large. The truth is merely that there are many kinds of seconds in the world,
although one can at the same time quite well say that a second is the time in
which light with its constant velocity $c$ runs the path $c$.

% --- p. 204 ---
This is no real contradiction either; but space and time have hitherto been
conceived so incorrectly that it looks like one. As soon as they are conceived
in their right relation to each other, everything comes into the most
beautiful order. ---

Thus does Einstein's doctrine issue. The contradictions in it are pushed
further and further out into the distance, as criminals were once sent to
Botany Bay. In the end it is the fundamental determinations space and time
themselves that suffer for it. Further than that one cannot go, unless one
were to take one very last step and maintain that contradiction and
non-contradiction are also at bottom one and the same.

Let us next look a little at Lorentz's doctrine.

One can, as has been indicated, quite well arrive at this doctrine by starting
from Einstein's and gradually removing all its contradictions and
exaggerations. Lorentz has however gone another way. Since he \emph{relied on
the theory of light} and found it necessary to assume a resting ether, he
explains Michelson's result to himself in the only way still possible, spoken
of before: by conjecturing that the motion has brought with it a contraction
of the apparatus in the direction of the translation. And as soon as the
kindred experiments mentioned by Dr.\ H. had come in, he sought further to
explain to himself the origin of the contraction in such a way that precisely
the results obtained must come about. Thereby arose his more elaborated
hypothesis, whose germs are however already indicated in the first. He puts
forward his assumption with great reserve, constantly calls it merely a
hypothesis, and states ,that one cannot yet know at all how generally the
% PRINTER'S DEFECT, transcribed as printed: »udtaler ,at man endnu slet ikke
% kan vide« -- the comma is set AFTER the word space, tight against »at«,
% where the sense wants »udtaler, at«. There is no matching mark anywhere in
% the sentence, so it is a misplaced comma and not an opening quotation mark.
% Reproduced in the same position here.
contraction in question will set in; about that only experience will gradually
be able to instruct us. His investigation has merely shown that such a
contraction \emph{would be able} to come about by the ways indicated without
giving contradiction; but it has not shown that it \emph{must} come about.
Neither does Lorentz set up any dogma about the velocity of light as the
greatest in the world; he says quite simply that his hypothesis holds only for
velocities smaller than that of light. He proceeds, in other words, constantly
in a \emph{correctly scientific} manner.

What now are the consequences of his assumptions?

First and foremost this, that everywhere where the conjectured contraction of
things on account of their motion really sets

% --- p. 205 ---
in --- and it is of course \emph{possible} that it will \emph{always} set in
--- there it will go with us as it went with Michelson: if we work in the
electrodynamic or optical domain in a uniformly rectilinearly moved system
without taking into account that it possibly moves, we shall on account of the
contraction apparently get the same results as we should have got if the
system had been at rest.\footnote{In the face of questions where time also
plays a part, this will however strictly speaking only occur provided that we
either ourselves set our clocks to an apparently correct rate by the aid of
light signals, or that the contraction itself will really bring it about that
they get this rate, even if the system has previously been at rest and their
rate then corresponded to rest. This question has hardly yet been
sufficiently investigated.}

Further, the transformation formulae between two such systems will of course,
as purely mathematical consequences, get the same \emph{form} in Lorentz as in
Einstein.

But they will get a widely different \emph{significance}.

Lorentz is not a relativist. He constantly says ``the Einsteinian principle of
relativity'' when he speaks of it. He himself holds absolute rest and
unaltered motion to be two different states, and the system at rest in
relation to the ether is for him the really resting one, the ``distinguished''
system which the relativists deny because we cannot for the present point it
out \emph{experimentally}. If $A$'s system is such a one, then for Lorentz
$A$'s time is the true time and $B$'s only an apparent or conditioned one, and
the corresponding holds of course of their measurements. There is therefore in
Lorentz no question of the \emph{contradiction} that \emph{both $A$ and $B$
are right}. \emph{With him truth is only one}, even if we cannot always find
out where it lies. We know in any case that at most only one of the two can be
right. Lorentz thus denies neither thought's demands nor its results, and he
constantly distinguishes clearly between rational truth and what is
experimentally attainable. Space, time and simultaneity therefore preserve
with him undisturbed their former character, and what in Einstein appeared as
the artificial and the paradoxical, precisely because it was conceived as
something absolute or solely valid, presents itself in Lorentz as the quite
intelligible and natural, being clearly stamped as the expression of the
\emph{composite} situation which arises partly from the contraction that plays
its part and partly from the deliberate disregard of the system's possible or
real motion.
% In the two longest runs on p. 205 the italic capitals A and B stand INSIDE
% the letterspaced run; they are set as mathematics within \emph, as in the
% Danish. Kroman is here answering the sentence of his own 1915 article that
% Høffding later cites, »Sandheden er ikke dobbelt. Sandheden er kun én«
% (Relativitetsprincipet, p. 15) -- rendered there "Truth is not double.
% Truth is only one."

\begin{center}
\rule{0.16\textwidth}{0.5pt}
\end{center}
% END OF THE ARTICLE. No signature, no dateline: the text simply stops after
% the footnote, and a short centred tail-piece follows -- a broken rule of
% graduated dashes, rendered as a plain short rule, as at the end of the 1915
% article.

\end{document}
